\documentclass[11pt,a4paper]{article}
\usepackage[T1]{fontenc}
\usepackage{lmodern}
\usepackage[margin=25mm]{geometry}
\usepackage{amsmath,amssymb,amsthm,booktabs,longtable,array}
\usepackage[hidelinks]{hyperref}
\usepackage{xurl}
\hypersetup{
    colorlinks=true,
    linkcolor=blue,
    urlcolor=blue,
    citecolor=blue,
    pdftitle={The Ninth-Dimensional Snake-in-the-Box Problem},
    pdfauthor={広中依緒璃}
}

\usepackage{microtype}
\newtheorem{theorem}{Theorem}[section]
\newtheorem{proposition}[theorem]{Proposition}
\newtheorem{lemma}[theorem]{Lemma}
\newtheorem{corollary}[theorem]{Corollary}
\newtheorem{conjecture}[theorem]{Conjecture}
\theoremstyle{definition}
\newtheorem{definition}[theorem]{Definition}
\theoremstyle{remark}
\newtheorem{remark}[theorem]{Remark}
\newcommand{\Q}{Q}
\newcommand{\Ham}{d_{\mathrm H}}
\title{The Ninth-Dimensional Snake-in-the-Box Problem:\\
A Sequence-Based Conjecture, Verified Bounds, and Construction Obstructions}
\author{Iori Hironaka}
\date{\today}
\begin{document}
\maketitle
\begin{abstract}
The maximum length of an induced path in the binary hypercube is known through dimension eight, with values $1,2,4,7,13,26,50,98$. Subtracting a common exponential baseline gives two constant three-term correction blocks. Extrapolating their doubling suggests the independent conjecture $a(9)=196$. A second equally exact fit predicts 195, demonstrating the limited force of this numerical evidence. This paper presents proofs and explicit finite checks for a verified 191-edge witness, the reduction of optimality to exclusion of the next length, a counterexample to an endpoint-maximality argument, a cycle-deletion construction, the exact two-layer joining criterion, a spectral bound $a(9)\le255$ obtained from Huang's degree theorem, and a square-face argument refining the bound to 254. The spectral bound, explicit witnesses, local singleton-square lemma, and square-face upper bound all have independent machine checks. The square-face proof counts incidences from the four possible corners and bounds selected two-edge corners by the number of internal path vertices. Restricted searches exclude several construction families, including all two-vertex deletion neighborhoods of 131 published 191-edge examples. These restricted solver results have not been independently certified. No 192-edge improvement, 196-edge construction, or matching upper bound has been established. The exact ninth-dimensional optimum remains open in this investigation.
\end{abstract}
\setcounter{tocdepth}{1}
\tableofcontents
\clearpage

\section{Problem, conventions, and prior observations}
\begin{definition}
The graph $\Q_n$ has vertex set $\{0,1\}^n$, with adjacency $u\sim v$ when $\Ham(u,v)=1$. A snake of length $L$ is a sequence of distinct vertices $v_0,\ldots,v_L$ such that
\[
v_i\sim v_j\quad\Longleftrightarrow\quad |i-j|=1\qquad(i\ne j).
\]
Length counts edges. Thus a length-$L$ snake has $L+1$ vertices. Let $a(n)$ denote the largest possible length.
\end{definition}
Finiteness of $\Q_n$ ensures that a maximum exists. A transition sequence $t_1\cdots t_L$, with $0\le t_i<n$, specifies a path from the zero vertex by
\[
v_0=0,\qquad v_i=v_{i-1}\mathbin{\oplus}2^{t_i}.
\]
Integers encode binary words; $\oplus$ is exclusive-or. Translation by exclusive-or with a fixed word, coordinate permutations, and reversal preserve the snake conditions.

The exact values through dimension eight are listed in OEIS A099155 \cite{oeis}. Carlson and Hougen constructed a length-98 snake in dimension eight; the upper bound was established by exhaustive search by Östergård and Pettersson \cite{ostergard}. The universal bound $a(8)=98$ is used as a cited result here; this project checks a witness of that length but does not reproduce the complete eight-dimensional classification.

Orland et al. supply 191-edge snakes in dimension nine and a public collection of 131 inequivalent examples \cite{orland,data}. Optimality is unresolved. The collection is a discovered bank, with no proved global completeness. Its examples cannot be extended at their endpoints; that observation alone gives no maximum-length theorem.

\paragraph{Evidence conventions.}
Theoretical statements are accompanied by proofs, and explicit witnesses are checked against the defining pairwise conditions. Search outcomes are recorded separately: \textsc{Infeasible} covers the specified model and restrictions, while \textsc{Unknown} establishes neither existence nor nonexistence. The restricted search exclusions have not been independently certified. The results include standard consequences and independently checked reductions; no new mathematical record is claimed.

\section{Independent hypothesis from the first eight terms}
\begin{center}
\begin{tabular}{rrrr}
\toprule
$n$ & $a(n)$ & $3\cdot2^{n-3}$ & correction\\
\midrule
3&4&3&1\\4&7&6&1\\5&13&12&1\\
6&26&24&2\\7&50&48&2\\8&98&96&2\\
\bottomrule
\end{tabular}
\end{center}
The baseline coefficient is also recovered from successive differences within the two triples:
\[
\frac{7-4}{2^4-2^3}=\frac{13-7}{2^5-2^4}
=\frac{50-26}{2^7-2^6}=\frac{98-50}{2^8-2^7}=\frac38.
\]
Continuing the correction as $1,1,1,2,2,2,4,4,4,\ldots$ gives
\begin{equation}\label{eq:model}
\widehat a(n)=3\cdot2^{n-3}+2^{\lfloor(n-3)/3\rfloor},\qquad n\ge3.
\end{equation}
The terms at dimensions one and two are assigned separately as 1 and 2.

\begin{conjecture}\label{conj:196}
The maximum snake length in $\Q_9$ is $196$ edges.
\end{conjecture}
This is a low-confidence conjecture. It is selected as a concrete target for construction and refutation, not inferred as a consequence of the known values. Equation~\eqref{eq:model} predicts $3\cdot64+4=196$.

\begin{proposition}[Model recurrence]
For $n\ge3$, the model satisfies
\[
\widehat a(n+3)+6\cdot2^{\lfloor(n-3)/3\rfloor}=8\widehat a(n).
\]
\end{proposition}
\begin{proof}
Put $t=n-3$ and $c=2^{\lfloor t/3\rfloor}$. Since $\lfloor(t+3)/3\rfloor=\lfloor t/3\rfloor+1$, the left side is
$24\cdot2^t+2c+6c=8(3\cdot2^t+c)$.
\end{proof}
At $n=6$ the recurrence predicts $8\cdot26-12=196$. This is the same evidence as the correction table, not an independent argument.

\begin{proposition}[The supplied terms do not determine the ninth prediction]
The model
\[
\widetilde a(n)=3\cdot2^{n-3}+1+\left\lfloor\frac{n-3}{3}\right\rfloor
\qquad(n\ge3)
\]
agrees with every supplied term, after assigning the first two terms separately, and gives $\widetilde a(9)=195\ne196$.
\end{proposition}
\begin{proof}
For $n=3,4,5$, both corrections equal 1. For $n=6,7,8$, both equal 2. At $n=9$, the corrections are respectively 3 and 4. Direct substitution proves the claim.
\end{proof}
Only two correction blocks have been observed. No construction currently explains why either proposed correction should describe the actual extremal function.

\section{Exact targets for an optimality proof}
\begin{lemma}[Prefix restriction]\label{lem:prefix}
Every prefix of a snake is a snake.
\end{lemma}
\begin{proof}
Restrict the vertex sequence to its initial indices. Distinctness is inherited. Adjacency between any two retained indices has the same value and the same consecutive-index condition as in the original sequence.
\end{proof}

\begin{theorem}[Next-length exclusion]\label{thm:next}
For every integer $b\ge0$, all snakes in $\Q_9$ have length at most $b$ if and only if no snake has length $b+1$.
\end{theorem}
\begin{proof}
The upper bound directly excludes $b+1$. Conversely, a snake of any length larger than $b$ has a prefix of length $b+1$ by Lemma~\ref{lem:prefix}.
\end{proof}

\begin{corollary}
Conjecture~\ref{conj:196} is equivalent to the conjunction of a 196-edge construction and nonexistence of a 197-edge snake. A 196-edge construction would contain a 192-edge prefix and would refute the proposed upper bound 191.
\end{corollary}
\begin{proof}
Apply Theorem~\ref{thm:next} with $b=196$, and restrict a 196-edge snake to its first 193 vertices.
\end{proof}
For the competing optimum 191, the missing exclusion is the finite statement that no sequence $v_0,\ldots,v_{192}\in\{0,1\}^9$ is distinct and satisfies
\[
\Ham(v_i,v_j)=1\quad\Longleftrightarrow\quad |i-j|=1.
\]
The pairwise condition is required for every $0\le i<j\le192$. Neither next-length exclusion has been proved.

\section{Witnesses, endpoint obstruction, and an invalid inference}
\begin{theorem}[Verified lower bound]\label{thm:lower}
A 191-edge snake exists in $\Q_9$.
\end{theorem}
\begin{proof}
The first transition sequence in Appendix~\ref{app:witness} is the published Appendix E example of \cite{orland}. Starting at zero reconstructs 192 vertices. For each $0\le i<j\le191$, direct calculation gives $v_i\ne v_j$ and $\Ham(v_i,v_j)=1$ exactly when $j=i+1$. This is a finite verification of $\binom{192}{2}=18\,336$ unordered pairs using the displayed transition sequence, and establishes a snake of length 191.
\end{proof}

\begin{definition}
A snake $P=(v_0,\ldots,v_L)$ is blocked at its right endpoint when every neighbor $w$ of $v_L$ either equals a vertex $v_i$ with $i<L$, or is adjacent to such a vertex.
\end{definition}
\begin{lemma}[Blocking prevents endpoint extension]\label{lem:block}
A right-blocked snake cannot be extended by appending one vertex while retaining its entire ordered sequence.
\end{lemma}
\begin{proof}
An appended $w$ must be adjacent to $v_L$. If $w=v_i$ for an earlier index, distinctness fails. If $w\sim v_i$ for $i<L$, the appended vertex has a chord to an index at least two positions earlier. Both alternatives contradict the snake definition.
\end{proof}

\begin{theorem}[Endpoint maximality does not imply maximum length]
There is a 93-edge snake in $\Q_9$ blocked at both endpoints, although a 191-edge snake exists in the same cube.
\end{theorem}
\begin{proof}
The second sequence in Appendix~\ref{app:witness} specifies 94 vertices. Direct calculation verifies distinctness and the induced adjacency condition for all $\binom{94}{2}=4\,371$ unordered pairs. The appendix gives an earlier blocking vertex for each of the nine endpoint moves, in both the original and reversed orderings. Lemma~\ref{lem:block} excludes extension at either end. Theorem~\ref{thm:lower} supplies a strictly longer snake.
\end{proof}

\begin{theorem}[A sufficient universal blocking condition]
If every 191-edge snake in $\Q_9$ is blocked at its right endpoint, then $a(9)=191$.
\end{theorem}
\begin{proof}
If a 192-edge snake existed, its 191-edge prefix would admit extension by its remaining final vertex. This contradicts the universal blocking premise. Apply Theorem~\ref{thm:next} and the verified lower bound.
\end{proof}
The premise concerns every 191-edge snake. Blocking the 131 recorded classes is insufficient until completeness of the collection is established. Repeated recovery of the same classes is not a completeness proof.

\section{Cycles as an alternative construction route}
\begin{theorem}[Deleting a cycle vertex]
Deleting one vertex from an induced cycle of $C$ vertices leaves a snake of length $C-2$.
\end{theorem}
\begin{proof}
Order the cycle as $v_0,\ldots,v_{C-1},v_0$ and delete $v_{C-1}$. Among the remaining vertices, the only surviving cycle edges join consecutive indices. The closing edge involved the deleted vertex; the induced-cycle hypothesis forbids all other edges. The resulting path has $C-1$ vertices and $C-2$ edges.
\end{proof}
\begin{corollary}
An induced cycle with 194 vertices in $\Q_9$ would disprove the upper bound 191.
\end{corollary}
\begin{proof}
Deletion gives a 192-edge snake.
\end{proof}
In particular, a 192-vertex coil gives only a 190-edge snake through this deletion operation. Edge and vertex counts cannot be interchanged.

\section{The two-layer construction and its obstruction}
Write a vertex of $\Q_9$ as $(v,b)$ with $v\in\Q_8$ and $b\in\{0,1\}$. Vertices in the same layer retain their original adjacency. Across the layers,
\[
(u,0)\sim(v,1)\quad\Longleftrightarrow\quad u=v.
\]
The same description defines the product of an arbitrary graph with a binary edge. Taylor discusses such lifts in \cite{taylor}.

\begin{theorem}[Exact joining condition]\label{thm:join}
Let $P=(p_0,\ldots,p_r)$ and $R=(q_0,\ldots,q_s)$ be induced paths in a graph. The ordered sequence
\[
(p_0,0),\ldots,(p_r,0),(q_0,1),\ldots,(q_s,1)
\]
is an induced path precisely when
\[
p_i=q_j\quad\Longleftrightarrow\quad i=r\ \hbox{and}\ j=0.
\]
When valid, its length is $r+s+1$.
\end{theorem}
\begin{proof}
The layer labels ensure that vertices in different layers are distinct. Each individual path retains its distinctness and induced adjacency conditions. Every cross-layer edge is caused by equality of the base vertices. In the displayed order there is exactly one consecutive pair across layers, at $(p_r,0),(q_0,1)$. Therefore the induced condition requires exactly the stated equality and forbids all other equalities. These conditions also suffice, since all possible pairs have been covered. The vertex count is $(r+1)+(s+1)$, hence the edge count is $r+s+1$.
\end{proof}

\begin{corollary}[Single-crossing length constraints]
Assume the known bound $a(8)=98$. A single-crossing snake in $\Q_9$ has length at most 197. Length 196 requires layer lengths $(98,97)$ or $(97,98)$; length 197 requires $(98,98)$.
\end{corollary}
\begin{proof}
The two layer segments are eight-dimensional snakes. Thus $r,s\le98$ and $L=r+s+1$. For $L=196$, $r+s=195$ forces the two possibilities; for $L=197$, both upper bounds must be attained.
\end{proof}
\begin{theorem}[Integer realization of the layer construction]
The map $J(u,b)=u+256b$, for $0\le u<256$ and $b\in\{0,1\}$, is injective and preserves adjacency in both directions between $\Q_8\mathbin{\square}K_2$ and the integer encoding of $\Q_9$. Consequently, compatible layer paths of lengths 98 and 97 give a 196-edge snake under the same definition as the verified witness. Compatible paths of lengths 98 and 98 give a 197-edge snake and would refute the conjectured upper bound 196.
\end{theorem}
\begin{proof}
The intervals $[0,255]$ and $[256,511]$ are disjoint, so equality of images forces equality of both components. Flipping a coordinate $k<8$ preserves the layer, because
\[
(u+256)\mathbin{\oplus}2^k=(u\mathbin{\oplus}2^k)+256,
\qquad u\mathbin{\oplus}2^k<256.
\]
Flipping coordinate 8 exchanges $u$ and $u+256$. These are all nine possible flips. Thus adjacency within a layer is exactly eight-dimensional adjacency, while adjacency between layers requires equal base vertices. Compose this map with Theorem~\ref{thm:join}. The lengths become $98+97+1=196$ and $98+98+1=197$. Conversely, every unwanted equality between base vertices creates a nonconsecutive cross-layer edge, so the compatibility condition remains necessary after applying $J$.
\end{proof}
This argument covers snakes that cross a chosen coordinate exactly once. A general snake can cross every coordinate repeatedly. Its intersection with a layer can then be a disjoint union of paths, so bounding that entire intersection by a single-snake length would be invalid.

The length-98 seed in Appendix~\ref{app:witness} is taken from \cite{records} and independently checked against all pairwise conditions. Two tests were performed. First, truncate either endpoint of an isometric copy to length 97, choose both joining endpoints, and enumerate all $8!$ coordinate permutations. Translation is forced by identifying the chosen endpoints. The $2\cdot2\cdot2\cdot8!=322\,560$ placements, counted with possible duplication, have minimum overlap 15. Theorem~\ref{thm:join} requires overlap one. Second, a constraint solver searched for an arbitrary length-97 partner avoiding every seed vertex except the chosen endpoint. Both endpoint instances were infeasible. These restricted computational exclusions have not been independently certified; other length-98 seeds and multiple-crossing constructions remain possible.

\section{A global upper bound}
\begin{theorem}[Huang's degree theorem]\label{thm:huang}
Let $n\ge1$. If $H\subseteq\{0,1\}^n$ and $|H|>2^{n-1}$, then the subgraph induced by $H$ has a vertex of degree at least $\sqrt n$.
\end{theorem}
\begin{proof}[Proof argument]
Define real symmetric matrices recursively by $A_0=(0)$ and
\[
A_n=\begin{pmatrix}A_{n-1}&I\\I&-A_{n-1}\end{pmatrix}.
\]
Induction gives $A_n^2=nI$. Its nonzero entries are precisely signed hypercube edges. For $n\ge1$, its trace is zero, so its eigenvalues $+\sqrt n$ and $-\sqrt n$ each have multiplicity $2^{n-1}$. The subspace of vectors supported on $H$ has dimension $|H|$. If $|H|>2^{n-1}$, it has nonzero intersection with the positive eigenspace, by the dimension inequality for intersections. Choose a nonzero vector $y$ in that intersection and $q\in H$ maximizing $|y_q|>0$. Then
\[
\sqrt n\,|y_q|=|(A_ny)_q|
\le\sum_{p\in H,\ p\sim q}|y_p|
\le \deg_H(q)|y_q|.
\]
Division proves the assertion. This is Huang's argument \cite{huang}, included to make the dependence of the upper bound explicit.
\end{proof}

\begin{lemma}[Compatibility of the integer and bit-function encodings]
Map $v\in\{0,\ldots,511\}$ to its nine binary digits $E(v)(k)$. This map is injective. If $E(u)$ and $E(v)$ differ in exactly one coordinate, then $v=u\oplus2^k$ for some $k<9$.
\end{lemma}
\begin{proof}
All digits at indices at least nine vanish because $v<2^9$. Equality of the first nine digits therefore gives equality at every digit and hence equality of integers. If exactly digit $k$ differs, exclusive-or with $2^k$ changes that digit and preserves every other digit. Equality of all digits then gives the stated integer identity.
\end{proof}

\begin{lemma}[Degree bound for an encoded snake]
The set of bit-function encodings of a snake's vertices induces maximum degree at most two.
\end{lemma}
\begin{proof}
For a vertex with index $i$, any encoded neighbor corresponds, by the preceding lemma, to an integer-encoded adjacent vertex. The snake condition forces its index to be $i-1$ or $i+1$, when those indices exist. There are at most two such vertices. Distinctness and injectivity of $E$ preserve the vertex count.
\end{proof}

\begin{theorem}[Verified universal bound]\label{thm:upper}
Every snake in $\Q_9$ has at most 255 edges. Consequently,
\[
\boxed{191\le a(9)\le255.}
\]
\end{theorem}
\begin{proof}
A snake of at least 256 edges would have at least 257 vertices, exceeding $2^8$. Theorem~\ref{thm:huang} would give an induced degree of at least $\sqrt9=3$, contradicting the preceding degree bound. Combine the upper bound with Theorem~\ref{thm:lower}.
\end{proof}
The bound is a standard consequence of Huang's theorem. Its contribution here is a checked global guarantee and an explicit connection between the two cube representations. It does not distinguish 191, 195, or 196.


\section{Square faces and the role of endpoints}
For a vertex set $S\subseteq\Q_9$ and $x\in S$, let $\alpha_S(x)$ be the number of two-dimensional coordinate faces whose intersection with $S$ is exactly $\{x\}$. Such a face will be called a singleton square at $x$. Square-face counting occurs in the cycle bounds reviewed by Lukito and van Zanten \cite{lukito}. Their cycle argument uses a lower bound of two singleton squares at each selected vertex. A direct extension of that condition to every internal vertex of every open snake is false.

\begin{theorem}[An internal-vertex counterexample]\label{thm:square-counterexample}
There is a 137-edge snake $P=(v_0,\ldots,v_{137})$ in $\Q_9$ with $v_{127}=0$ and $\alpha_{V(P)}(0)=1$. The vertex $v_{127}$ is at least ten path edges from either endpoint.
\end{theorem}
\begin{proof}
Use the transition sequence in Appendix~\ref{app:square-counterexample}, starting at the integer vertex 145. Direct calculation checks all $\binom{138}{2}=9\,453$ unordered vertex pairs: the vertices are distinct, and adjacency holds precisely for consecutive indices. It also gives
\[
v_0=145,\qquad v_{127}=0,\qquad v_{137}=288.
\]
The selected neighbors of zero are 1 and 2. Every square at zero involving coordinate 0 or 1 therefore contains another selected vertex. For $2\le a<b\le8$, neither $2^a$ nor $2^b$ is selected, and the opposite vertex $2^a+2^b$ is selected except when $(a,b)=(2,3)$. Thus the unique singleton square is $\{0,4,8,12\}$. Finally, the path distances from $v_{127}$ to the endpoints are 127 and 10.
\end{proof}
The Hamming distances from zero to the two endpoints are 3 and 2. Distance along the path is therefore insufficient for ignoring endpoint effects in this local cube configuration. This example concerns a 137-edge path; it makes no claim about the corresponding local condition for all longest paths.

\begin{lemma}[At least one singleton square]\label{lem:one-square}
If $S$ is the vertex set of an induced path in $\Q_9$, then $\alpha_S(x)\ge1$ for every $x\in S$.
\end{lemma}
\begin{proof}
Translate $x$ to zero. Translation preserves adjacency and the induced-path property. Addition below is coordinatewise modulo two, and $e_i$ denotes the unit vector in coordinate $i$.

First suppose that the path has at most two edges. The three coordinate squares using the pairs $(0,1)$, $(2,3)$, and $(4,5)$ have mutually disjoint sets of nonzero vertices. If none were a singleton square, they would contribute three distinct selected vertices besides zero, exceeding the number of vertices in the path. Thus at least one is a singleton square.

Now suppose that the path has at least three edges, and assume for contradiction that no singleton square exists at zero. For any two distinct coordinates $i,j$ whose unit vectors are absent, the opposite vertex $e_i+e_j$ must then be selected.

Every vertex of a path with at least three edges has an internal path vertex as a neighbor. Choose such a neighbor $e_a$ of zero, and let $c$ be its other path neighbor. Write $c=e_a+e_i$, where $i\ne a$ because $c\ne0$. The unit vector $e_i$ is absent: otherwise $0,e_a,c,e_i,0$ would be a selected four-cycle, which an induced path cannot contain.

At most two unit vectors are selected, since zero has induced degree at most two. Choose a coordinate $b\notin\{a,i\}$ so that every selected unit vector is indexed by $a$ or $b$. If there is a second selected unit vector, take its coordinate; otherwise any coordinate outside $\{a,i\}$ suffices. Set
\[
U=\{0,\ldots,8\}\setminus\{a,b\}.
\]
Then $i\in U$, every unit vector indexed by $U$ is absent, and every pair vertex $e_j+e_k$ with distinct $j,k\in U$ is selected. Every triple vertex $e_j+e_k+e_l$ with distinct $j,k,l\in U$ is absent, because its three pair-vertex neighbors are selected and a path has maximum degree two.

Choose four distinct coordinates $j\in U\setminus\{i\}$; six are available. For each of these four coordinates, define
\[
p_j=e_i+e_j,\qquad q_j=e_a+e_i+e_j.
\]
All four $p_j$ are selected. The four $q_j$ are distinct neighbors of $c$, distinct also from its already selected neighbor $e_a$. Therefore at most one $q_j$ is selected, and at least three are absent.

For each absent $q_j$, the corresponding $p_j$ has at most one selected neighbor. Flipping coordinate $i$ or $j$ gives an absent unit vector; flipping any other coordinate in $U$ gives an excluded triple; and flipping $a$ gives the absent $q_j$. Only the neighbor $e_b+e_i+e_j$ can remain selected. Hence at least three distinct selected vertices have induced degree at most one. In a path with at least three edges, precisely the two endpoints have degree at most one. This contradiction proves the lemma.
\end{proof}

\begin{theorem}[Square-face refinement]\label{thm:square-upper}
Every snake in $\Q_9$ has at most 254 edges. Together with the explicit witness, this gives
\[
\boxed{191\le a(9)\le254.}
\]
\end{theorem}
\begin{proof}
A path with one vertex satisfies the assertion, so let $S$ be a path vertex set of size $v\ge2$. For $0\le a<b\le8$ and $x\in\Q_9$, write
\[
C_{ab}(x)=\{x,x+e_a,x+e_b,x+e_a+e_b\},
\]
where addition is coordinatewise modulo two. Define three indicators, each taking values in $\{0,1\}$:
\begin{align*}
s_{ab}(x)&=1\quad\Longleftrightarrow\quad C_{ab}(x)\cap S=\{x\},\\
w_{ab}(x)&=1\quad\Longleftrightarrow\quad x,x+e_a,x+e_b\in S,\\
e_{ab}(x)&=1\quad\Longleftrightarrow\quad C_{ab}(x)\cap S=\varnothing.
\end{align*}
All three indicators are zero when their stated condition fails. Let
\[
T=\sum_{a<b}\sum_{x\in\Q_9}s_{ab}(x),\qquad
W=\sum_{a<b}\sum_{x\in\Q_9}w_{ab}(x),\qquad
E=\sum_{a<b}\sum_{x\in\Q_9}e_{ab}(x).
\]
No square has all four vertices selected, since that would be a four-cycle in the induced path. For each square, considering its possible numbers of selected vertices, zero through three, gives
\begin{equation}\label{eq:local-square-identity}
2e_{ab}(x)+\sum_{y\in C_{ab}(x)}s_{ab}(y)
 +\sum_{y\in C_{ab}(x)}\mathbf{1}_S(y)
=2+\sum_{y\in C_{ab}(x)}w_{ab}(y).
\end{equation}
For three selected vertices, exactly one is adjacent to the other two, so the last sum is one. For zero, one, or two selected vertices, the last sum is zero; the other terms give the asserted identity directly.

For each fixed coordinate pair, the four translations taking $x$ to the corners of $C_{ab}(x)$ are permutations of $\Q_9$. Summing~\eqref{eq:local-square-identity} over all 512 choices of $x$ and all 36 coordinate pairs therefore gives
\begin{equation}\label{eq:anchored-square-count}
2E+4T+144v=36\,864+4W.
\end{equation}
Here every selected vertex contributes four times for each coordinate pair to the vertex-incidence sum.

Lemma~\ref{lem:one-square} implies $T\ge v$. A nonzero contribution to $W$ requires two distinct selected neighbors of its selected center, so the center is an internal path vertex. Since an induced path has degree at most two, at most one unordered coordinate pair contributes at any such center. There are $v-2$ internal vertices, hence $W\le v-2$. Using $E\ge0$ in~\eqref{eq:anchored-square-count},
\[
4v+144v\le36\,864+4(v-2),
\qquad\text{so}\qquad144v\le36\,856<144\cdot256.
\]
Thus $v\le255$, and the number of edges is $v-1\le254$.
\end{proof}
\paragraph{Verification status.}
The explicit counterexample, Lemma~\ref{lem:one-square}, the local identity~\eqref{eq:local-square-identity}, its global summation~\eqref{eq:anchored-square-count}, and the upper bound 254 are independently machine-checked. The upper bound is derived from the same induced-path definition used for the 191-edge witness; no census equations or optimality assumptions are supplied as premises. The earlier spectral bound 255 provides an independent, weaker check. No claim of a new best published upper bound is made. Neither a 196-edge construction nor a matching upper bound has been obtained.

\subsection{A restriction on one-singleton configurations}
A possible route to a stronger aggregate bound is to show that vertices with only one singleton square must lie near an endpoint. The following result narrows the local configuration without establishing that claim.

\begin{theorem}[The two continuations use different missing coordinates]\label{thm:split-arms}
Let $S$ be an induced path vertex set in $\Q_9$ containing zero. Suppose its selected unit vectors are exactly $e_a,e_b$, where $a\ne b$, and put $U=\{0,\ldots,8\}\setminus\{a,b\}$. Let $u,v\in U$ be distinct. Assume every pair vertex $e_s+e_t$ with distinct $s,t\in U$, other than possibly $e_u+e_v$, is selected, and every selected pair vertex indexed by $U$ is internal in the path. Suppose also that $e_a+e_i$ and $e_b+e_j$ are selected for some $i,j\in U$. Then
\[
(i,j)=(u,v)\quad\text{or}\quad(i,j)=(v,u).
\]
\end{theorem}
\begin{proof}
First suppose that $i\notin\{u,v\}$. Choose two distinct coordinates $k$ outside $\{a,b,i,u,v\}$; at least four are available. For either choice, $e_i+e_k$ is selected and internal. Its only possible selected neighbors are $e_a+e_i+e_k$ and $e_b+e_i+e_k$. Indeed, deleting $i$ or $k$ gives an absent unit vector. Adding any other coordinate $t\in U$ gives a triple with three selected pair neighbors: none of its three pairs can be the exceptional pair $\{u,v\}$, since both $i$ and $k$ lie outside that pair. Such a triple is forbidden by degree at most two.

Both possible neighbors of the internal vertex $e_i+e_k$ must therefore be selected. The two choices of $k$ give two distinct neighbors of $e_a+e_i$, in addition to its selected neighbor $e_a$, a contradiction. Hence $i\in\{u,v\}$. Interchanging $a,b$ proves $j\in\{u,v\}$.

Suppose next that $i=j=u$. Choose three distinct $k\in U\setminus\{u,v\}$. Each selected internal vertex $e_u+e_k$ has possible selected neighbors only among
\[
e_a+e_u+e_k,\qquad e_b+e_u+e_k,\qquad e_u+e_v+e_k.
\]
All other triples have three selected pair neighbors, while unit neighbors are absent. At least one of the first two displayed vertices must therefore be selected, since an internal vertex requires two selected neighbors. However, $e_a+e_u$ already has selected neighbor $e_a$, and can have at most one of the three vertices $e_a+e_u+e_k$ as another selected neighbor. The same argument applies to $e_b+e_u$. Thus at most two of the three choices of $k$ can be covered, a contradiction. Interchanging $u,v$ excludes $i=j=v$. The two remaining ordered pairs are exactly those asserted.
\end{proof}

The theorem has an independent machine check under its stated geometric premises. If zero has exactly one singleton square and all endpoints have Hamming distance at least three from zero, the missing square supplies $u,v$, and every selected pair vertex is internal. The root and its two unit neighbors are internal as well; their second neighbors supply the two continuations. Thus the theorem identifies the remaining local arrangement. It does not prove that such a full path is impossible.

\begin{proposition}[Radius-three constraints do not suffice]\label{prop:radius-three}
There is an induced linear forest on 60 vertices of $\Q_9$, consisting of 14 pairwise separated paths, such that every selected vertex lies within Hamming distance three of zero, every selected vertex at distance at most two has degree two, and zero has exactly one singleton square.
\end{proposition}
\begin{proof}
Take the 14 ordered paths in Appendix~\ref{app:radius-three}. Direct checks show that every path is induced, their vertex sets are disjoint, and no edge joins different paths. Counting selected neighbors gives degree two at every selected vertex of weight at most two. Every listed vertex has weight at most three. Among the 36 coordinate squares through zero, exactly the square using coordinates two and three contains no other listed vertex. These finite checks have independent machine checks.
\end{proof}

The forest satisfies the local degree and acyclicity requirements but is disconnected. It is therefore not a counterexample to a statement about a single full snake. Its existence shows that these radius-three requirements alone cannot complete the proposed endpoint-proximity argument. Connections through higher-weight vertices, or another global restriction, remain necessary.


\subsection{Cycle constraints beyond four-cycles}
\begin{lemma}[A finite subset has a vertex of degree at most one]\label{lem:subset-low-degree}
Let $P=(p_0,\ldots,p_m)$ be an induced path and let $A$ be a nonempty subset of its vertex set. Some vertex of $A$ has at most one neighbor in $A$.
\end{lemma}
\begin{proof}
Choose the least index $i$ for which $p_i\in A$. Any neighbor $p_j\in A$ of $p_i$ satisfies $j\ge i$ by minimality and $|j-i|=1$ by the induced-path condition. Thus $j=i+1$. There is at most one such vertex.
\end{proof}
Consequently, a finite nonempty vertex set having two distinct neighbors within the set at every vertex cannot be contained in an induced path. This excludes cycles of arbitrary length and also applies to disconnected unions of cycles. The lemma and the validity of all 41 recorded cycle cuts in the radius-five experiment have been independently machine-checked.

\section{Computational investigations and their exact scope}
\subsection{Published bank audit and endpoint regrowth}
Independent Python checks verified all 131 published nine-dimensional examples, including their transition and vertex lists, every pairwise adjacency condition, and blocked endpoints. The total number of unordered pairs checked was $131\binom{192}{2}=2\,402\,016$. Canonical transition sequences, normalized by first coordinate occurrence and compared with their reversals, distinguish 131 classes under translation, coordinate permutation, and reversal. This audit makes no completeness claim.

Endpoint regrowth retained a prefix of each example in each orientation, then searched every legal continuation to 193 total vertices. A vertex is forbidden after it has been used, or when it is adjacent to a retained vertex preceding the current endpoint. Adding a vertex blocks all neighbors of the previous endpoint. The resulting depth-first search covers the permitted continuations. Node limits produce \textsc{Unknown} rather than a nonexistence result.
\begin{center}\begin{tabular}{rrrrr}
\toprule
Removed & Retained & Excluded & Unresolved & Nodes\\\midrule
32&160&262&0&488,450\\
40&152&262&0&16,561,030\\
48&144&209&53&66,968,134\\\bottomrule
\end{tabular}\end{center}
The final row used one million nodes per orientation. These exhaustive-search results concern only the listed prefixes and have not been independently certified.

\subsection{Constraint models and vertex replacement}
The exact model selects cube vertices and directed cube-edge arcs. A single directed circuit through an auxiliary vertex encodes one connected path; deleting the auxiliary vertex gives its endpoints. Unselected cube vertices use self-loops. Every cube edge whose two endpoints are selected is required to appear in the path. Together these constraints enforce distinctness, connectivity, and absence of chords. Small feasible and infeasible thresholds in dimensions three through five were reproduced. The unrestricted 192-edge run returned \textsc{Unknown} after 300 seconds.

A broader local model permits both endpoints to change and imposes only
\[
|V(P)\setminus V(T)|\le k,\qquad |V(T)|=193,
\]
where $P$ is a recorded 191-edge snake. Thus changes may occur anywhere in the path. For $k=2$, all 131 models were infeasible, with about 241 seconds total solver time. In vertex-set terms, a 192-edge improvement would have to delete at least three vertices from each tested representative. A two-vertex deletion allows three added vertices, since the target has one extra vertex overall. This excludes a restricted neighborhood, without excluding a distant snake. The corresponding model with $k=8$ returned \textsc{Unknown} on the first three seeds within 30 seconds each. The implementations and exclusions have not been independently certified.

\subsection{Other searches and unresolved cases}
Four structured searches for a 194-vertex coil used quotient-cube constructions with extra generators of weights 2, 4, 6, and 8. All were unresolved at their time limits; a small control lifted to a valid 14-vertex coil. No coil improvement follows.

For the fixed eight-dimensional seed, searches for partner lengths 93, 92, and 91 at both endpoints were each unresolved after 30 seconds. A relaxation imposing only induced degree at most two admits 106 vertices in the permitted region, so that local degree condition alone is insufficient. Retaining 50 initial seed vertices in either orientation and exhaustively completing to length 98 recovered only the original completion, after 1,926,810 nodes in total. Retaining 42 vertices reached its twenty-million-node limit in both orientations. Neither result classifies all length-98 snakes.

\subsection{Joint search with both layer paths variable}
The fixed-seed obstruction does not exclude other pairs. A joint model therefore selects both projected paths, requires each to start at zero, and forbids every shared nonzero projected vertex. A circuit constraint with one auxiliary vertex enforces a single connected path in each layer; selected adjacent cube vertices must be consecutive in that path. The only cross-layer edge is then the intended join at zero.

The two first coordinate directions must differ, since a common first neighbor would be an additional shared vertex. A common coordinate permutation may therefore send those directions to coordinates zero and one. With and without this symmetry restriction, all 42 tested cases in dimensions two and three agreed with independent exhaustive enumeration.

\begin{proposition}[A checked starting pair]\label{prop:joint-seed}
There are induced paths of lengths 96 and 94 in $\Q_8$ whose vertex sets intersect only at their common initial vertex. Their two-layer join has 191 edges.
\end{proposition}
\begin{proof}
Use the two transition sequences in Appendix~\ref{app:joint-seed}, both starting at zero. They reconstruct 97 and 95 vertices. All $\binom{97}{2}=4\,656$ and $\binom{95}{2}=4\,465$ within-path pairs satisfy the distinctness and induced-adjacency conditions. Checking all $97\cdot95=9\,215$ cross pairs gives equality exactly for the two initial vertices. Reversing the first path and applying the two-layer joining criterion gives $96+94+1=191$ edges. These finite checks and the joining deduction have independent machine checks.
\end{proof}

The pair was extracted from the first example in the supplied bank: coordinate five changes only once, at edge index 96, when edge indices start at zero. Deleting that coordinate and applying a common translation and coordinate permutation gives the displayed pair. This is a decomposition of a published example, not a new record. Fixing both selected vertex sets to this pair gave a feasible solver control.

Three searches then left both paths variable. Each used four workers, random seed 13, and a 120-second limit:
\begin{center}
\begin{tabular}{rrrrl}
\toprule
First length & Second length & Joined target & Starting hint & Status\\
\midrule
98 & 97 & 196 & none & \texttt{UNKNOWN}\\
98 & 98 & 197 & none & \texttt{UNKNOWN}\\
96 & 95 & 192 & 96,94 pair & \texttt{UNKNOWN}\\
\bottomrule
\end{tabular}
\end{center}
The hint in the last run does not fix either path. These timeouts establish no existence or impossibility result. The model and its search implementation have not been independently certified; the explicit starting pair and the mathematical joining criterion have been checked separately.


\subsection{A radius-five exclusion and its verification boundary}
Fix the normalized radius-two selected set
\[
F=\{0,1,2,5,10\}\cup
\{2^a+2^b:2\le a<b\le8,\ (a,b)\ne(2,3)\},
\]
where integers encode binary vertices. In the ball of Hamming radius $r$ about zero, the model selects exactly $F$ at weights at most two. Every selected vertex strictly inside the ball has degree one or two; selected vertices of weight at most two have degree two. Boundary vertices have degree at most two. The number of interior vertices of degree one is bounded by a prescribed budget. Connectedness is not imposed. Cycle cuts exclude specified cyclic components.

At radius four, a forest exists with interior endpoint budget two. At radius five, the solver reports infeasibility for budgets two and three, and finds a forest with budget four. Allowing cycles restores feasibility at budget two. The feasible assignments have been checked independently. These findings distinguish an obstruction involving both endpoint count and cycles from the weaker radius-three conditions in Proposition~\ref{prop:radius-three}.

A Boolean encoding of the radius-five, budget-two model records the exact cycle cuts. Removing irrelevant clauses and replacing auxiliary counters by direct restrictions on triples of endpoint indicators gives a 567-clause core involving 252 Boolean variables. Its certificate consists of 398 derived clauses obtained by reverse unit propagation: assume the negation of the proposed clause, propagate the cited clauses, and obtain a contradiction. The resulting Boolean contradiction has been independently machine-checked through ordinary logical proof steps. This certifies the explicit formula; no geometric exclusion is asserted without the remaining interpretation proof.

The remaining verification boundary is explicit. The interpretation of each input clause as a necessary consequence of the geometric model must be checked. Generic degree and endpoint clause implications, and all recorded cycle cuts, have been checked separately. The reduction from an arbitrary one-singleton vertex to the normalized selection $F$ also remains to be completed. Until these steps are established, the experiment does not prove endpoint proximity or improve the interval $191\le a(9)\le254$.

\subsection{A quantitative target for the endpoint argument}
\begin{proposition}[Conditional improvement from a singleton surplus]
With the quantities in~\eqref{eq:anchored-square-count}, suppose a snake with $v\ge2$ vertices satisfies $T+92\ge2v$. Then it has at most 250 edges.
\end{proposition}
\begin{proof}
The checked incidence identity, $E\ge0$, and $W\le v-2$ give
\[
140v+4T\le36\,856.
\]
Substituting $T\ge2v-92$ yields $148v\le37\,224$. Since $148\cdot252=37\,296$, it follows that $v\le251$ and hence $v-1\le250$. The arithmetic implication has been independently machine-checked; its new aggregate premise remains unproved.
\end{proof}
A Hamming ball of radius two in $\Q_9$ contains $1+9+\binom92=46$ vertices. Thus, if every vertex with exactly one singleton square were proved to lie within distance two of an endpoint, at most 92 vertices could be exceptional. Lemma~\ref{lem:one-square} would then suggest the required aggregate inequality: nonexceptional vertices contribute at least two singleton incidences and exceptional vertices at least one. The endpoint statement and its checked aggregation remain open. The bound 250 is therefore a conditional target, not an established global bound.

\section{Scope of the results and remaining obligations}
The preceding arguments establish explicit lower-bound witnesses, reductions of the optimality question, conditional constructions, and a universal upper bound. The numerical extrapolation remains conjectural. The single-crossing length constraints use the cited eight-dimensional maximum $a(8)=98$; this paper does not repeat the exhaustive proof of that result. No compatible pair yielding length 196 or 197 has been found. The restricted search results cover precisely the recorded seeds and construction families, and provide no global exclusion.

The remaining obligations for Conjecture~\ref{conj:196} are exact: construct and check a 196-edge snake, then exclude length 197. A 192-edge witness would already refute the competing optimum 191. Conversely, excluding every 192-edge snake would establish 191 and refute Conjecture~\ref{conj:196}. At present neither outcome has been achieved. The numerical hypothesis is retained provisionally while both directions remain under investigation.

\appendix
\section{Explicit transition sequences and blockers}\label{app:witness}
All coordinates are numbered from zero. Spaces only separate groups; start at the zero vertex unless another initial vertex is stated, and apply each coordinate flip in order.
\subsection{Verified 191-edge nine-dimensional witness}
\begin{quote}\small\ttfamily
0120314506 2106502631 4321035023 6210623607 4063260126\\
3206483026 3153263120 3623402362 0162360750 6326012632\\
0652362013 4120132610 2156210263 0210674601 2063201260\\
4816201351 2013261021 4621026302 1067560120 6
\end{quote}
Its last vertex is 118. For moves in coordinates $0,\ldots,8$, respectively, earlier blocking indices are
\[
(19,37,13,13,27,19,189,43,91).
\]
For each coordinate $k$, put $w=v_{191}\mathbin{\oplus}2^k$. Direct calculation gives $w=v_i$ or $\Ham(w,v_i)=1$ for the corresponding displayed index $i<191$.

\subsection{Verified 93-edge endpoint-maximal counterexample}
\begin{quote}\small\ttfamily
8624301806 7812851450 7860736287 5286237138 0746235243\\
8145783485 6180720370 4716078615 6326120174 036
\end{quote}
The terminal vertex is 289. Blocking indices for its nine moves are
\[
(1,33,51,91,45,1,91,81,45).
\]
For the reversed ordering, whose terminal vertex is zero, the blocking indices in that reversed ordering are
\[
(25,19,43,27,19,61,21,23,91).
\]
The reversed ordering is $v_{93},v_{92},\ldots,v_0$. Its pairwise distances are unchanged. Each listed blocker precedes the reversed endpoint and either equals the candidate next vertex or has Hamming distance one from it.

\subsection{Verified 98-edge eight-dimensional seed}
\begin{quote}\small\ttfamily
0123104154 3146340134 6743104274 0147514016 3104154314\\
2340134274 3104674014 7561340142 7410431641 04764013
\end{quote}
The seed is attributed to Carlson and Hougen by the source page \cite{records}. The exact transition and vertex data for all four witnesses accompany the paper in JSON files.

\subsection{A 137-edge counterexample for singleton squares}\label{app:square-counterexample}
Start at the integer vertex 145, rather than zero, and use the following coordinate transitions:
\begin{quote}\small\ttfamily
0132518601 5823563751 0231051201 8510231061 2714831024\\
7201352471 3814601234 6187510231 0512018510 2376320158\\
1021501320 1871023158 1321531023 1581320
\end{quote}
The resulting ordered path has 138 distinct vertices. Its vertex at index 127 is zero, and its last vertex is 288. At zero, the only singleton coordinate square uses coordinates 2 and 3.

\subsection{Compatible paths of lengths 96 and 94}\label{app:joint-seed}
Both projected paths start at zero in $\Q_8$. For the first path use
\begin{quote}\small\ttfamily
0213047013 0412031052 1031243103 2672301324 1301237103\\
1240541240 1342147314 2130214236 0135231421 032142
\end{quote}
For the second path use
\begin{quote}\small\ttfamily
1031240741 2401342145 3142130214 2367324123 0124132714\\
2130450130 4120310721 0312431032 6523013241 3012
\end{quote}
The projected vertex sets have intersection $\{0\}$ only. Reverse the first sequence of vertices, place it in layer zero, cross at zero, and follow the second path in layer one.

\subsection{A radius-three linear forest}\label{app:radius-three}
The following lists give vertices as decimal integers encoding nine-bit words. Each row is one path; the rows are distinct connected components of the induced graph on their union.
\begin{align*}
P_{1}&=(134,132,140,136,138,10,2,0,1,5,261,260,268,264,265)\\
P_{2}&=(22,20,28,24,25)\\
P_{3}&=(38,36,44,40,41)\\
P_{4}&=(49,48,50)\\
P_{5}&=(70,68,76,72,73)\\
P_{6}&=(81,80,82)\\
P_{7}&=(97,96,98)\\
P_{8}&=(145,144,146)\\
P_{9}&=(161,160,162)\\
P_{10}&=(193,192,194)\\
P_{11}&=(273,272,274)\\
P_{12}&=(289,288,290)\\
P_{13}&=(321,320,322)\\
P_{14}&=(385,384,386)
\end{align*}
The union has 60 vertices, including zero. It is a linear forest with 14 components, not a single snake.

\begin{thebibliography}{9}\raggedright
\bibitem{oeis} OEIS Foundation. \emph{A099155: Maximum length of a simple path with no chords in the $n$-dimensional hypercube}. \url{https://oeis.org/A099155}. Accessed 13 October 2026.
\bibitem{ostergard} P.~R.~J. Östergård and V.~H. Pettersson. Exhaustive Search for Snake-in-the-Box Codes. \emph{Graphs and Combinatorics} 31 (2015), 1019--1028. \url{https://doi.org/10.1007/s00373-014-1423-3}.
\bibitem{orland} P.~Orland et al. \emph{New Snake-in-the-Box Records via Snakepit Surgery and Learned Construction}. arXiv:2607.15270v3, 30 September 2026. \url{https://arxiv.org/html/2607.15270v3}.
\bibitem{data} Math-AI-Caltech. \emph{Snake-in-the-Box: public record data}. \url{https://github.com/Math-AI-Caltech/Snake-in-the-Box}. Downloaded archive SHA-256 is recorded in the accompanying bank audit.
\bibitem{records} T.~Ace. \emph{Snake-in-the-Box lower bounds and transition sequences}. \url{https://www.minortriad.com/snake/}. Accessed 13 October 2026.
\bibitem{taylor} T.~Taylor. \emph{Constructing longer snakes and improved asymptotic bounds in hypercubes}. arXiv:2609.27483v1, 23 September 2026. \url{https://arxiv.org/html/2609.27483v1}.
\bibitem{huang} H.~Huang. Induced subgraphs of hypercubes and a proof of the Sensitivity Conjecture. \emph{Annals of Mathematics} 190 (2019), 949--955. \url{https://annals.math.princeton.edu/2019/190-3/p06}.
\bibitem{lukito} A.~Lukito and A.~J. van Zanten. A New Non-Asymptotic Upper Bound for Snake-in-the-Box Codes. \emph{Journal of Combinatorial Mathematics and Combinatorial Computing} 39 (2001), 147--156. \url{https://combinatorialpress.com/article/jcmcc/Volume%20039/vol-039-paper%2012.pdf}.
\end{thebibliography}
\end{document}
